\chapter{基于 Attention-MAPPO 的双层协同优化方法}
\label{cha:method}

\section{整体框架}
\label{sec:framework}

本文框架由两个协同 MARL 决策层组成：

\begin{itemize}
    \item \textbf{上层 attention-MAPPO}：每架 UAV 作为一个 agent，根据自身位置、邻居 UAV 状态、覆盖 UE 和请求负载信息输出连续轨迹动作（2D 方向向量）。
    \item \textbf{下层 constrained attention offload MAPPO}：每架 UAV 对其当前覆盖的请求 slot 输出离散卸载动作；实际动作空间包含 local、MBS 和指定协作 UAV 三类动作。
\end{itemize}

在每个环境 step 中，系统按如下顺序执行：
\begin{enumerate}
    \item 环境生成当前 UAV、UE、链路和请求状态
    \item 下层读取请求级观测与动作 mask，输出卸载动作
    \item 环境执行服务请求处理，统计 latency、energy、DSR、MBS load ratio
    \item 上层输出轨迹动作，环境更新下一时刻位置
\end{enumerate}

\begin{figure}[!htbp]
    \centering
    \includegraphics[width=0.85\textwidth]{docs/figures/fig_system_workflow.png}
    \bicaption{单时隙系统工作流程}{Single time-slot system workflow}
    \label{fig:system-workflow}
\end{figure}

\begin{figure}[!htbp]
    \centering
    \includegraphics[width=0.85\textwidth]{docs/figures/fig_hmarl_framework.png}
    \bicaption{双层 MARL 框架}{Hierarchical MARL framework}
    \label{fig:hmarl-framework}
\end{figure}

\section{上层 attention-MAPPO 轨迹控制}
\label{sec:upper-layer}

上层将每架 UAV 作为一个 agent。第 $i$ 架 UAV 的观测由自身状态、邻居 UAV 状态和覆盖 UE 状态三部分组成：
\begin{equation}
    o_i^U = [s_i^{self}, S_i^{nbr}, S_i^{ue}]
    \label{eq:upper-obs}
\end{equation}

其中 $s_i^{self}=[p_i,c_i]$ 包含 UAV 归一化二维位置 $p_i$ 和缓存向量 $c_i$；$S_i^{nbr}=\{\Delta p_{ij}|j\in\mathcal{N}_i\}$ 包含邻居 UAV 相对位置；$S_i^{ue}=\{\Delta p_{ik}, r_k, b_k|k\in\mathcal{C}_i\}$ 包含覆盖 UE 的相对位置、请求特征和电池状态。

具体地，首先将自身状态、邻居状态和 UE 状态分别编码为隐表示：
\begin{equation}
    e_i=f_{self}(s_i^{self}),\quad h_{ij}^{nbr}=f_{nbr}(s_{ij}^{nbr}),\quad h_{ik}^{ue}=f_{ue}(s_{ik}^{ue})
    \label{eq:encodings}
\end{equation}

随后对邻居 UAV 集合和覆盖 UE 集合分别计算 scaled dot-product attention。以邻居 UAV 集合为例：
\begin{equation}
    Q_i=W_Qe_i,\quad K_{ij}=W_Kh_{ij}^{nbr},\quad V_{ij}=W_Vh_{ij}^{nbr}
    \label{eq:attention-qkv}
\end{equation}
\begin{equation}
    \alpha_{ij}^{nbr}=\text{softmax}_j\left(\frac{Q_iK_{ij}^{T}}{\sqrt{d}}\right),\quad c_i^{nbr}=\sum_{j\in\mathcal{N}_i}\alpha_{ij}^{nbr}V_{ij}
    \label{eq:attention-output}
\end{equation}

对 UE 集合同理得到 $c_i^{ue}$。最终将自身表示、邻居上下文和 UE 上下文拼接后得到策略特征：
\begin{equation}
    z_i=\text{MLP}([e_i,c_i^{nbr},c_i^{ue}])
    \label{eq:policy-feature}
\end{equation}

上层策略根据 $z_i$ 输出 2D 方向向量，转换为飞行距离：
\begin{equation}
    d_{moved} = \text{clip}(||a||, 0, 1) \cdot v_{max} \cdot \tau
    \label{eq:flight-distance}
\end{equation}

上层奖励为第 \ref{sec:optimization-objective} 节定义的系统级综合奖励 $R_t$，所有 UAV 共享同一基础奖励值。训练采用 CTDE（Centralized Training Decentralized Execution）框架。

\begin{figure}[!htbp]
    \centering
    \includegraphics[width=0.85\textwidth]{docs/figures/fig_trajectory_before_after.png}
    \bicaption{轨迹前后对比图}{Trajectory before-after comparison}
    \label{fig:trajectory-compare}
\end{figure}

\begin{figure}[!htbp]
    \centering
    \includegraphics[width=0.85\textwidth]{docs/figures/fig_service_coverage_fairness.png}
    \bicaption{覆盖公平性热力图}{Service coverage fairness heatmap}
    \label{fig:fairness-heatmap}
\end{figure}

\section{下层 constrained attention offload MAPPO}
\label{sec:lower-layer}

下层同样采用多智能体建模。每架 UAV 是一个 lower-layer agent，处理其当前关联的 bounded request slots（最多 30 个）。每个有效 request slot 的离散动作空间如表 \ref{tab:action-space} 所示。

\begin{table}[!htbp]
    \centering
    \bicaption{下层离散动作空间}{Lower-layer discrete action space}
    \label{tab:action-space}
    \begin{tabular}{|c|l|l|}
        \hline
        动作编号 & 动作含义 & 说明 \\
        \hline
        0 & Local UAV execution & 在当前 UAV 本地执行 \\
        1 & MBS execution & 通过回传链路卸载到宏基站 \\
        $2+i$ & Cooperative UAV $i$ execution & 选择编号为 $i$ 的邻居 UAV 执行，$i=0,\ldots,N-1$ \\
        \hline
    \end{tabular}
\end{table}

因此当前 $N=5$ 时，下层 actor 的实际动作数为 $|A|=2+N=7$。

下层观测同样按 UAV 构造。第 $i$ 架 UAV 的下层观测为：
\begin{equation}
    o_i^L=[g_i^{self},G_i^{req}]
    \label{eq:lower-obs}
\end{equation}

其中 $g_i^{self}$ 包含 UAV 归一化位置、当前服务请求队列比例、邻居数量比例和 UAV-MBS 回传速率；$G_i^{req}$ 由最多 30 个 request slot 组成。每个有效请求 slot 包含请求大小、请求编号、deadline、priority、本地缓存命中、协作可用性、本地队列长度、local/cooperative/MBS 时延相对 deadline 的比例、最佳邻居计算份额以及 UE-UAV 相对位置等特征。

下层 actor 使用 request attention 编码同一 UAV 内多个请求之间的相对重要性：先将 UAV 自身特征和每个请求特征编码到隐空间，再对 request slots 做多头自注意力，最后对每个 request slot 输出 local、MBS 或 cooperative UAV 动作的 logits。critic 采用集中式信息估计 value。下层奖励函数为：
\begin{equation}
    R_{lower} = w_{succ} \cdot DSR + w_{coop} \cdot C_{ratio} - w_{dead} \cdot (1 - DSR) - w_{lat} \cdot \frac{L}{M \cdot T_{penalty}} - w_{en} \cdot \frac{E}{N \cdot E_{ref}^{offload}} - w_{mbs} \cdot M_{ratio} - \lambda_{dsr} \cdot \max(0, \tau_{dsr} - DSR) - \lambda_{mbs} \cdot \max(0, M_{ratio} - \tau_{mbs})
    \label{eq:lower-reward}
\end{equation}

其中权重为：$w_{succ}=0.4$, $w_{coop}=0.08$, $w_{dead}=3.0$, $w_{lat}=4.0$, $w_{en}=1.2$, $w_{mbs}=0.35$。

\section{质量感知动作 Mask}
\label{sec:action-mask}

请求级动作 mask 用于减少明显无效的探索：

\textbf{空请求 slot mask}：若某个 slot 没有真实服务请求，则该 slot 不产生真实卸载动作。此保护在所有实验中始终开启。

\textbf{质量感知 mask}（quality mask）：
\begin{itemize}
    \item \textbf{Cooperative mask}：当系统不存在可用邻居、协作路径时延过高（$> 1.5 \times$ deadline）、协作相对非协作路径无明显优势（$> 1.35 \times$ 本地时延）或协作计算份额过低（$< 25\%$）时，屏蔽 cooperative 动作。
    \item \textbf{MBS mask}：当 MBS 路径时延显著超过 deadline clip 阈值时，屏蔽 MBS 动作。
\end{itemize}

\begin{figure}[!htbp]
    \centering
    \includegraphics[width=0.85\textwidth]{docs/figures/fig_action_mask_decision.png}
    \bicaption{动作 mask 与下层决策示意图}{Action mask and lower-layer decision illustration}
    \label{fig:action-mask}
\end{figure}

\section{Lagrange 约束机制}
\label{sec:lagrange}

为显式控制 DSR 与 MBS 负载之间的权衡，本文引入 Lagrange 乘子 $\lambda_{dsr}$ 和 $\lambda_{mbs}$：
\begin{equation}
    P_t = \lambda_{dsr} \cdot \max(0, \tau_{dsr} - DSR_t) + \lambda_{mbs} \cdot \max(0, M_t - \tau_{mbs})
    \label{eq:lagrange-penalty}
\end{equation}

约束目标为 $\tau_{dsr} = 0.18$，$\tau_{mbs} = 0.03$。$\lambda$ 在每个 rollout 后更新：
\begin{equation}
    \lambda_{dsr} \leftarrow \text{clip}(\lambda_{dsr} + \eta_{\lambda} \cdot \bar{v}_{dsr}, 0, \lambda_{max})
    \label{eq:lambda-dsr-update}
\end{equation}
\begin{equation}
    \lambda_{mbs} \leftarrow \text{clip}(\lambda_{mbs} + \eta_{\lambda} \cdot \bar{v}_{mbs}, 0, \lambda_{max})
    \label{eq:lambda-mbs-update}
\end{equation}

其中 $\eta_{\lambda} = 0.1$ 为 Lagrange 学习率，$\lambda_{max} = 10.0$ 为上界，$\bar{v}$ 为 rollout 平均约束违反量。

\section{奖励函数设计}
\label{sec:reward-design}

本文统一采用线性归一化奖励函数。系统级综合奖励 $R_t$（第 \ref{sec:optimization-objective} 节）直接使用归一化后的指标线性加权，下层局部训练奖励 $R_{lower}$（第 \ref{sec:lower-layer} 节）同样采用线性归一化形式。

归一化分母的设计如下：
\begin{itemize}
    \item \textbf{Latency 归一化}：以 $M \cdot T_{penalty}$ 为分母（$M=100$ 个 UE，$T_{penalty}=20$s 为未服务惩罚时延），因此归一化时延 $\bar{L} \in [0, 1]$ 表示当前总时延占最差情况的比例。
    \item \textbf{Energy 归一化}：系统级奖励使用 $N \cdot E_{ref}$ 为分母（$N=5$ 架 UAV，$E_{ref}=5000$J 为每架 UAV 单步参考能耗）；下层奖励使用 $N \cdot E_{ref}^{offload}=125000$J（$N=5$, $E_{ref}^{offload}=25000$J 为下层每架 UAV 单步最大参考能耗）。
\end{itemize}

与早期版本使用的对数归一化相比，线性归一化的优势在于：(i) 各指标贡献在数值上可加可比，便于权重调参；(ii) 不同策略间的 reward 值可以在统一口径下直接对比；(iii) 避免了 log 函数在接近零值时的数值不稳定性。

\section{复杂度分析}
\label{sec:complexity}

设 UAV 数量为 $N=5$，每架 UAV 的最大请求 slot 数为 $K=30$，下层实际动作数为 $|A|=2+N=7$。若直接构建端到端联合 MAPPO，需要同时处理 $N$ 个连续轨迹动作和最多 $N \times K=150$ 个请求级离散卸载动作，单步离散卸载组合可达 $|A|^{NK}$，并带来严重的样本效率与 credit assignment 问题。本文采用双层分解后，上层只处理 $N$ 个 UAV 的轨迹动作（连续 2D），下层只在每个 UAV 的 $K$ 个请求 slot 上进行局部离散决策。attention 编码带来的主要额外开销为 $O(N^2)$ 或 $O(K^2)$。

\section{本章小结}
\label{sec:summary-chap04}

本章详细描述了基于 Attention-MAPPO 的双层协同优化方法。首先介绍了整体框架设计，包括上层轨迹控制模块和下层卸载决策模块；然后详细介绍了上层 attention-MAPPO 的网络结构和训练流程；接着介绍了下层 constrained attention offload MAPPO 的设计，包括质量感知动作掩码和拉格朗日约束机制；随后介绍了奖励函数设计和复杂度分析。该框架通过双层分解降低了直接联合动作建模的复杂度，并通过 attention 机制建模了 UAV 间的协同关系。
